% Construcción de la innovación de la distribución implícita.
\begin{tikzpicture}[
  n/.style={caja=#1,text width=3.1cm,minimum height=13mm},
  node distance=6mm and 7mm]
\node[n=qazul] (obs) at (0,0) {\textbf{Cambio observado}\\ $\Delta x_t(u)=x_t(u)-x_{t-1}(u)$};
\node[n=qnaranja] (base) at (0,-2.3) {\textbf{Información base} $\mathcal{I}^{base}_t$\\ retorno ya observado,
  volatilidad, mercado y sector, gap, calendario};
\node[n=tenue] (mod) at (4.1,-2.3) {\textbf{Modelo de expectativa}\\ ajustado solo con fechas\\ anteriores a $t$};
\node[n=tenue] (esp) at (8.2,-2.3) {\textbf{Parte esperada}\\ $\widehat{\E}\bigl[\Delta x_t(u)\mid\mathcal{I}^{base}_t\bigr]$};
\node[circle,draw=tinta2,line width=0.7pt,minimum size=6mm,font=\sffamily\bfseries] (resta) at (8.2,0) {$-$};
\node[n=qazul] (inn) at (12.2,0) {\textbf{Innovación}\\ $\epsilon_t(u)$};
\node[n=qnaranja] (uso) at (12.2,-2.3) {\textbf{Variable candidata}\\ para pronosticar $\P$\\ después de $t$};
\draw[flecha] (obs) -- (resta);
\draw[flecha] (base) -- (mod);
\draw[flecha] (mod) -- (esp);
\draw[flecha] (esp) -- (resta);
\draw[flecha] (resta) -- (inn);
\draw[flecha] (inn) -- (uso);
\node[nota,anchor=north] at (6.1,-3.25) {La innovación no prueba información privada ni causalidad: es un candidato cuyo valor
  predictivo se mide fuera de muestra.};
\end{tikzpicture}
